\scanpage{100}
Considérons, d'ailleurs deux torseurs $\unsure{\gamma}$-trivialisés $\deleted{\unsure{P}}P_1$, $P_2$ sur les mod. \uline{inv.} $L_1,L_2$, l'homomorphisme can
\diagram{\begin{tikzcd}[column sep=large]
E_1\otimes E_2 \arrow[r,no head,"x_1\otimes x_2\mapsto x_1*x_2"] \arrow[d,phantom,"\wr"] & E_1*E_2=E \arrow[d,phantom,"\wr"]\\
(\mathcal A_1\otimes L_1)\otimes(\mathcal A_2\otimes L_2) & \mathcal A\otimes L
\end{tikzcd}}
avec $L=L_1\otimes L_2$.
On en déduit
\[
\mathcal A_1\otimes\mathcal A_2\xrightarrow{\tau_A}\mathcal A
\]
qu'on va aussi dénoter
\[
\xi_1\otimes\xi_2\longmapsto\xi_1*\xi_2
\]
et qu'on voudrait étudier. On a d'autre part, dans le cas $\chi=2$ et quand $\unsure{P_1},P_2$ correspondent à des \deleted{$\mathcal A\longrightarrow$} rev. quadratiques $X_1,X_2$ de $X$, alors $E$ à un rev. quad. de \unsure{type}
\[
\mathcal X=X_1\wedge X_2
\]
de $X$, à cause de
\[
X_1\times X_2\longrightarrow\mathcal X
\]
induit par $(*)$, un hom ``transposé''
\[
\mathcal A\xrightarrow{i_A}\mathcal A_1\otimes\mathcal A_2
\]
mais c'est simplement le transposé de
\[
E_1\otimes E_2\xrightarrow[\inserted{x_1\otimes x_2\mapsto x_1*x_2}]{\tau_E}E
\]
qui a \inserted{donc} un sens sans utiliser les structures de rev. i.e. les formes quadratiques $Q_1,Q_2$, ni les conditions $\chi=2$.
\mspage{99}
